\chapter{General Introduction}

This chapter introduces the internship context, the problem the host company
asked to have solved, the objectives set for the one-month engagement, and
the organization of the rest of this report.

\section{Internship Context}

This report covers a one-month internship carried out at ST2i, a Tunisian IT
services and consulting company (a subsidiary of the STUDI Group), between
July 6 and July 31, 2026. The internship was organized as four one-week
Scrum sprints, each ending with a working, demoable increment of the
delivered system rather than a single delivery at the very end.

\section{Problem Statement}

ST2i bids on public and international tenders that require staffing one or
more \textit{expert positions} defined in a Terms of Reference (TDR)
document: each position specifies a mandatory diploma, a minimum number of
years of experience, required certifications, languages, and sector
experience. Before submitting a bid, ST2i must determine, for every required
position, which of its own employees are actually qualified.

Today this step is manual: a business analyst reads the TDR and tries to
recall or search which internal experts might fit. This is slow, and more
importantly it is error-prone in a way that has direct financial
consequences -- missing a genuinely qualified expert can mean passing on a
winnable bid, while missing a disqualifying gap can mean committing to staff
who do not actually meet a client's mandatory requirements.

\section{Objectives}

The internship's objective was to design and build \textbf{AdéquIA}, a
decision-support system that automates this matching step: given a parsed
TDR and the internal CV database, produce for each required position a
ranked list of candidates with a compatibility score, a breakdown by
criterion, flagged anomalies, and a citation-backed rationale that a human
reviewer can approve, reject, or override.

Because the internship's timeline is one month, scope was deliberately kept
to the matching, scoring, anomaly-detection, and reporting pipeline itself,
plus the reviewer-facing dashboard needed to actually use it -- tender
discovery and bid submission logistics were explicitly out of scope, as was
any model fine-tuning or self-hosted inference infrastructure.

\section{Report Organization}

Chapter~\ref{chap:host-company} presents the host company and the business
environment the project sits in. Chapter~\ref{chap:project-management}
covers the Agile/Scrum methodology adopted for the internship, including the
product backlog and the four-sprint plan. Chapters~4 onward each cover one
sprint in turn -- sprint goal, backlog, analysis, design, implementation,
testing, review, and retrospective -- following the same structure so
progress across sprints is directly comparable. A general conclusion closes
the report.
